\documentclass[12pt]{article}
\usepackage[utf8]{inputenc}
\usepackage[T1]{fontenc}
\usepackage{hyperref}
\title{Ndzwani Comorian Language Resources for K-12 Curriculum Development}
\author{Compiled by Claude Sonnet 5.5 (claude-sonnet-5-5) for Andr\'as Kornai}
\date{}
\begin{document}
\maketitle

\section{Introduction}
Ndzwani Comorian (Shindzuani, Shinzwani, Anjouanais; ISO 639-3 wni) is spoken on Ndzuani (Anjouan), Comoros, and belongs to the eastern group of Comorian languages together with Shimaore of Mayotte. It was historically written in Arabic-based Ajami, and the French colonial administration introduced Latin script.

\section{Primary Educational Resources}
I found no school primers or textbooks in Shindzuani in the search results.

\section{Linguistic and Reference Materials}
The main dictionary found is Harriet Ottenheimer's \emph{Comorian-English/English-Comorian (Shinzwani) Dictionary} (AG Press, Manhattan, KS, 2011, 403 pp.): \url{https://glottolog.org/resource/reference/id/551282}. Ottenheimer's 2001 article on dictionary construction and orthographic choice, ``Spelling Shinzwani'': \url{https://glottolog.org/resource/reference/id/149583}. The missionary William Elliott compiled a vocabulary of about 900 words with 98 sample sentences, plus a grammar, in the 1820s. This is mentioned in the overview at \url{https://en.wikipedia.org/wiki/Comorian_languages}. I did not find a modern standalone grammar.

\section{Religious and Community Resources}
None specific to Ndzwani were found.

\section{Digital and Audio Resources}
None found.

\section{Curriculum Development Context}
A bilingual dictionary and published orthographic discussion give a base for materials, but I found no existing curriculum. French and Arabic are used in Comorian schooling, and any curriculum would need local educators' input on orthography.

\end{document}